% Generated draft: requires mathematical and source review.
\documentclass[11pt]{article}
\usepackage[margin=25mm]{geometry}
\usepackage{fontspec}
\setmainfont{Latin Modern Roman}
\usepackage{amsmath,amssymb,mathtools,unicode-math}
\setmathfont{Latin Modern Math}
\usepackage{enumitem,needspace,xurl,hyperref}
\setlength{\parindent}{0pt}
\setlength{\parskip}{5pt}
\setlength{\emergencystretch}{3em}
\newcommand{\N}{\mathbb{N}}
\newcommand{\Z}{\mathbb{Z}}
\newcommand{\Q}{\mathbb{Q}}
\newcommand{\R}{\mathbb{R}}
\newcommand{\C}{\mathbb{C}}
\newcommand{\F}{\mathbb{F}}
\newcommand{\A}{\mathbb{A}}
\newcommand{\PP}{\mathbb{P}}
\newcommand{\E}{\mathbb{E}}
\newcommand{\eps}{\varepsilon}
\newcommand{\dd}{\,\mathrm{d}}
\begin{document}
\section*{UnsolvedMath ID: 701722}
\textbf{Source problem code:} catalogue entry
\subsection{Definitions and mathematical statement}
Use the abstract (discrete) group interpretation: the supplied record specifies no topology. A group \(G\) is amenable if there exists a finitely additive function \(\mu:\mathcal P(G)\to[0,1]\) such that \(\mu(G)=1\) and \(\mu(gA)=\mu(A)\) for every \(g\in G\) and \(A\subseteq G\). It is non-amenable if no such function exists.

A uniformly bounded representation of a group \(H\) on a complex Hilbert space \(\mathcal H\) is a homomorphism \(\pi:H\to\operatorname{GL}(\mathcal H)\), where \(\operatorname{GL}(\mathcal H)\) denotes the bounded invertible linear operators, such that
\[
\sup_{h\in H}\|\pi(h)\|<\infty.
\]
The group \(H\) is unitarisable if, for every complex Hilbert space \(\mathcal H\) and every uniformly bounded representation \(\pi\) on \(\mathcal H\), there is a bounded invertible operator \(S\) such that \(S\pi(h)S^{-1}\) is unitary for every \(h\in H\). It is non-unitarisable if there exists a uniformly bounded representation for which no such \(S\) exists.

For a positive integer \(n\), let \(G^n=\prod_{j=1}^{n}G\). Interpret the source notation \(G^{\infty}\) as the countable restricted direct product
\[
G^{\infty}=\bigoplus_{j=1}^{\infty}G
=\left\{(g_j)_{j\geq1}\in\prod_{j=1}^{\infty}G:\ g_j=1_G\text{ for all but finitely many }j\right\},
\]
with coordinatewise multiplication. The source explicitly specifies a direct sum (restricted product), but does not explicitly state its indexing set; countable indexing is the interpretation used here.

The questions proposed by N. Monod are:
\begin{enumerate}
\item[(a)] For every non-amenable group \(G\), does there exist an integer \(n\geq1\), allowed to depend on \(G\), such that \(G^n\) is non-unitarisable?
\item[(b)] As the weaker requested conclusion, is \(G^{\infty}\) non-unitarisable for every non-amenable group \(G\)?
\end{enumerate}
The supplied registry labels this problem open; this is a supplied status label, not an independently verified mathematical or literature assessment.
\subsection{Short English statement}
N. Monod asks whether every non-amenable group has a non-unitarisable finite direct power. The source also asks for the weaker conclusion that its restricted direct sum of infinitely many copies is non-unitarisable. This draft uses discrete groups, complex Hilbert spaces, and countably many copies in the infinite sum; the source leaves these conventions implicit.
\subsection{Sources}
\begin{itemize}
\item Supplied local UnsolvedMath registry record: permanent ID \(701722\), exact source problem code \texttt{catalogue entry}, title \emph{Kourovka Notebook Problem 20.64}; upstream revision \texttt{d371ee917ac9647fbb06cd6fc529deec6f3dda82}. The record was created and updated on \texttt{2026-10-16T14:00:00Z} and carries the status \texttt{open}.
\item \emph{The Kourovka Notebook (Unsolved Problems in Group Theory)}, Problem 20.64, proposed by N. Monod in 2022. The supplied background identifies the 20th issue (2022) and the 21st issue edition (2026), arXiv:1401.0300v48. The supplied collection description names E. I. Khukhro and V. D. Mazurov as editors of the collection covering Issues 1--21 (1965--2026). These bibliographic details are supplied attribution and have not been independently checked.
\end{itemize}
\end{document}
